% !TeX root = ../Book5.tex
\begin{figure}[htbp]
\centering
\begin{tikzpicture}[x=2.2cm,y=2.2cm,every node/.style={font=\small}]
  \fill[AlgebraBlue!12] (0,0)--(1.5,0)--(1,1)--(0,1.5)--cycle;
  \draw[->] (0,0)--(2.4,0) node[right] {\(a\)};
  \draw[->] (0,0)--(0,2.4) node[above] {\(b\)};
  \draw[AlgebraBlue] (0,1.5)--(1,1)--(1.5,0);
  \draw[dashed,gray] (.45,2.1)--(1.5,0);
  \draw[dashed,gray] (0,1.5)--(2.1,.45);
  \node[anchor=west] at (1.45,.6) {\(2a+b=3\)};
  \node[anchor=west] at (.55,1.6) {\(a+2b=3\)};
  \fill (0,0) circle[radius=1.5pt] node[below left] {\(0\)};
  \fill (1,1) circle[radius=1.5pt] node[above right] {\(\rho\)};
  \draw[fill=white] (1,0) circle[radius=1.5pt] node[below] {\(\omega_1\)};
  \draw[fill=white] (0,1) circle[radius=1.5pt] node[left] {\(\omega_2\)};
  \node[below] at (1.5,0) {\(3/2\)};
  \node[left] at (0,1.5) {\(3/2\)};
\end{tikzpicture}
\caption[\(A_2\) 中支配候选的有界交]{以基本权系数 \((a,b)\) 为坐标，阴影表示支配室与 \(\rho-\sum_i\mathbb R_{\ge0}\alpha_i\) 的交。
其中四个整数点只有两个黑点满足根格陪集条件；两个空心点不满足。}
\label{b5:fig:a2-dominant-bounded-intersection}
\end{figure}
